\documentclass[11pt,a4paper]{article}

\usepackage[margin=2.4cm]{geometry}
\usepackage[T1]{fontenc}
\usepackage{lmodern}
\usepackage{booktabs}
\usepackage{array}
\usepackage{amsmath}
\usepackage{enumitem}
\usepackage[hidelinks]{hyperref}
\usepackage{titlesec}

\setlength{\parskip}{0.55em}
\setlength{\parindent}{0pt}
\titlespacing*{\section}{0pt}{1.4em}{0.5em}
\titleformat{\section}{\large\bfseries}{\thesection.}{0.6em}{}

% a reviewer's point, run in ahead of our reply
\newcommand{\point}[1]{\textbf{#1}\ }

\title{\vspace{-1.4cm}\textbf{Response to reviewers}}
\author{}
\date{}

\begin{document}
\maketitle
\vspace{-2.2cm}

\begin{center}
\textit{Channel-Aware Backdoor Attacks Against Federated Infostealer Malware Classification\\
Using Dynamic API-Call and Network Representations} \quad (IWBIS)
\end{center}

\vspace{0.6em}

We thank all three reviewers. Every comment is answered below, with the section of the revised paper
that carries the change. Where we did not make a requested change, we say so and give the reason
rather than leaving it out.

Two of the reviewers asked for things that turned out to change our own reading of the results. The
significance tests Reviewers B and C requested showed that the red and green triggers are not
distinguishable from their controls at all, which is a stronger statement than the ``less effective''
we had written. Reviewer C's question about whether the blue channel's effectiveness comes from its
information content turned out to be answerable from the dataset construction code: the blue channel
is a deterministic function of the other two and carries no independent information. Both findings
are now in the paper.

\section{Table and figure numbering has changed}

Reviewers refer to tables by number, so please note the mapping before reading the responses. The
single-seed defense screening now precedes the multi-seed results, and the old Tables V and VII have
been merged.

\begin{center}
\begin{tabular}{@{}p{7.4cm}p{7.4cm}@{}}
\toprule
\textbf{In the submitted version} & \textbf{In this version} \\
\midrule
Table I, dataset summary & Table I, unchanged \\
Table II, experimental environment & Table II, unchanged (condensed from seven rows to four) \\
Table III, backbone selection & Table III, unchanged \\
Table IV, channel-aware backdoor sweep & Table IV, unchanged (caption now states seed 42) \\
Table V, main multi-seed results & \textbf{Table VI} \\
Table VI, single-seed defense screening & \textbf{Table V} \\
Table VII, multi-seed Multi-Krum results & \textbf{merged into Table VI} \\
\bottomrule
\end{tabular}
\end{center}

The merge removed two rows that Table VII duplicated from Table V exactly, the two
\emph{Backdoor, FedAvg} rows. No result was dropped. Figure numbering is unchanged: Figs.~1 to 7
still refer to the same figures.

Where a reviewer comment cites a table by its old number, our reply below gives the new number.

\section{Reviewer A}

\point{Originality should be supported by a clear comparison with closely related recent studies.}
Section~II-C now positions the work against backdoor attacks on malware
classifiers specifically: explanation-guided feature-space triggers (Severi et al., USENIX Security
2021) and family-selective triggers (Yang et al., IEEE S\&P 2023). Both attack static features of a
centralized model, whereas we attack a federated model through the channels of a dynamic-behavior
image. The abstract now also states the novelty in one sentence: the channels of the representation
are treated as the attack surface, and a trigger confined to the fusion channel alone matches one
spanning all three.

\point{The experimental procedure is not described with sufficient clarity; the stages are not
explicitly defined.} Section~III opens with a new paragraph that names each block of Fig.~1 in the
order the subsections take them. The procedure itself is now stated rather than implied.
Section~III-E gives the attack as implemented: the poison rate is 20\% of the attacker's 14
AgentTesla training images, that is two images per round, resampled every round; the trigger is a
white $12 \times 12$ square written at the bottom-right corner after normalization; the malicious
client is client~0 and participates in every round. Section~III-F gives the defense parameters,
including $f = 1$ and $|\mathcal{S}| = 2$ of five clients for Multi-Krum, and defines the trigger
control. Section~III-H states that ResNet18 is trained from scratch with AdamW, that the reported
model is the final-round global model with no best-validation selection, and that the per-round
curves are test-set curves.

\point{The workflow figure is insufficiently integrated; it should be referenced and its blocks
explained.} Fig.~1 is now referenced in the first sentence of Section~III, and the new opening
paragraph walks its blocks in order.

\point{Figure 2 is mentioned in II.B while it illustrates Dataset Creation, described in III.A.}
Fixed. The reference to Fig.~2 now sits in Section~III-A, next to the pipeline it depicts. Fig.~3,
which shows the processed representations, stays in III-B where the representations are described.

\point{Some tables are not mentioned anywhere in the paper.} All thirteen floats are now referenced
at least once in the text. Table~I is cited in III-A, Table~II in III-G, and Fig.~4 in III-D.

\point{The abstract should state the contribution or novelty more explicitly.} Added, as described
above.

\point{The discussion describes results but provides limited interpretation, and several claims lack
supporting references.} Section~IV now gives a measured mechanism for each of the two main results.
For the channel ordering, the mean intensity inside the trigger region across all 500 images is
241.7 in R, 81.5 in G and 10.2 in B, and 37.7\% of red trigger-region pixels are already at or above
254, where writing the trigger changes nothing. For Multi-Krum, the malicious client survives
selection in 8\% to 50\% of rounds across the six runs, and that rate correlates with the final ASR
(Pearson $r = 0.893$, $p = 0.017$, $n = 6$). The second result is connected to reports that robust
aggregation remains insufficient against stealthy objectives, citing [15] and [22]. The conclusion
positions the attack against FL backdoor work on general image tasks and against the color-space
backdoor.

\point{The reference list should represent recent studies related to the method, and the discussion
may need additional recent references.} We added the two 2021 and 2023 malware-backdoor references
named above and now use [15] and [22] in the discussion. We also removed two references that were
doing no work: one cited only for the cross-entropy loss, and one background citation in the opening
sentence whose claim the two remaining citations already support. The list is 25 entries, the same
count as the submitted version, with a more recent and more relevant composition. We verified
mechanically that every entry is cited and that the numbering follows the order of first citation.

\point{Standard of English: Fair.} The paper has been through two proofreading passes. Beyond
ordinary corrections, we unified the model notation, which was inconsistent between (1) and (3) to
(5), fixed a citation used to support a claim it contradicts, and corrected a statement that the
images are normalized to $[0, 255]$ when the code normalizes them to $[-1, 1]$.

\section{Reviewer B}

\point{The author should conduct significance testing.} Section~IV now reports Fisher exact tests on
the pooled triggered samples. Blue/fusion and full RGB differ from their trigger controls at
$p = 3.5 \times 10^{-19}$ and $p = 2.6 \times 10^{-18}$, and blue differs from red and from green at
$p = 4.0 \times 10^{-8}$. Multi-Krum reduces ASR against FedAvg at $p = 1.5 \times 10^{-8}$ for
blue/fusion and $p = 9.1 \times 10^{-7}$ for full RGB. Pooling across seeds is legitimate here
because the split is re-drawn for every seed, so the pooled samples are distinct, and the paper says
so where it reports the $p$-values. We also ran paired tests on clean performance across seeds: none
is significant, including the Multi-Krum utility drop, and the paper reports that with the caveat
that these tests are low-powered at $n = 3$.

This test changed one of our claims. We had written that the red and green triggers ``remain less
effective''. They are in fact indistinguishable from their own trigger controls, both at $p = 1$, so
the measured result is no detectable effect rather than a weaker one. Section~IV-B now says that.

\point{The author should explain the architecture of SmallCNN.} Section~III-C describes it: three
convolutional blocks of 32, 64 and 128 channels, each Conv $3 \times 3$ with batch normalization and
ReLU, max-pooling after the first two blocks, adaptive average pooling after the third, and a linear
layer to five classes. All backbones are trained from scratch.

\point{Tables I and II and Figures 1 and 4 exist but were never mentioned.} All four are now
referenced, as are the remaining floats.

\point{Keywords are too many.} Reduced from seven to five.

\point{The author should mention and discuss every table and figure.} Every float is now referenced
and discussed in at least one sentence.

\point{Contribution and originality: the contributions are empirical experiments.} We accept that
the contribution is empirical. We have made the specific claim sharper rather than broader: the
finding is that a trigger confined to the fusion channel alone is as effective as one spanning all
three channels, even though that channel is derived from the other two and carries no independent
information.

\section{Reviewer C}

\point{Backbone selection, the channel sweep and the single-seed defense screening all use seed 42,
so the selection may be biased toward one seed's behavior.} This is a fair criticism and the paper
now states it rather than leaving it implicit. The screening has been moved into its own subsection
(IV-C, Table~V) ahead of the multi-seed results, with the reason given: screening every defense
across every seed would multiply the training budget by the number of seeds for candidates that may
not be worth carrying forward, and all runs share one GPU. Section~IV-E then notes that the screening
uses seed 42, which is precisely the seed at which blue/fusion with Multi-Krum fails, so a defense
selected on one seed can misstate its behavior on the others. The captions of Tables III, IV and V
all now state that they use seed 42.

\point{Only three seeds; a deviation as large as $0.5111 \pm 0.4286$ suggests a bimodal outcome, not
partial mitigation.} You are right, and this was the most useful correction we received. Table~VI now
reports Multi-Krum per seed alongside the mean, with a note on both Multi-Krum rows stating that the
across-seed ASR range is at least 0.5, so the mean describes no individual run. The per-seed values
are 15/15, 3/15 and 5/15 for blue/fusion and 6/15, 15/15 and 6/15 for full RGB. Section~IV-E and the
conclusion now describe the outcome as bimodal rather than partial, and the abstract no longer quotes
the average ASR.

\point{The attack scenario tests only one source-target pair, so the generalization claim reaches
beyond the evidence.} We have not added a second pair, and the conclusion now lists this as a
limitation rather than arguing around it. The abstract qualifies the claim as holding ``in this
controlled IID setting''.

\point{The FL setting is IID-only although the motivation for FL is inherently non-IID.} Also not
addressed by new experiments, and also now stated as a limitation. We note that the split files
already carry a non-IID client assignment that these experiments do not use, so the experiment is
available to future work rather than blocked.

\point{No statistical significance testing is reported.} Added, as described under Reviewer B.

\point{The attacker's capability is underspecified; it is unclear whether model-replacement or
scaling as in [12] is employed.} Section~III-E now states that the malicious client's capability is
pure data poisoning: it trains and submits like any other client, with no update scaling and no model
replacement, so [12] is background rather than the attack implemented here.

\point{Multi-Krum on full RGB drops clean accuracy from 0.8400 to 0.7600, a trade-off that is
underemphasized in the discussion and the abstract.} The drop is now in both. Section~IV-E states it
next to the ASR reduction, and the abstract reports it as part of the Multi-Krum result.

\point{The conclusion states more confidence than $n = 3$ supports, and ``partially suppresses'' is
imprecise.} The phrase is gone. The conclusion now says that Multi-Krum reduced ASR bimodally rather
than partially, leaving the backdoor fully effective at 15/15 on one of the three seeds for each
trigger.

\point{The conclusion does not address the alternative explanation that the blue channel's
effectiveness may stem from its information content, which remains untested.} This turned out to be
testable without new experiments, and the answer is that the alternative can be ruled out. The
dataset construction code shows that the blue channel is the edge map of the average of the other
two, $B = \mathrm{FindEdges}((R+G)/2)$, which we verified against the shipped dataset by recomputing
it with zero difference on 400 of 400 images. The blue channel is therefore a deterministic function
of red and green and carries no independent information, so its effectiveness cannot come from
information the other channels lack. Section~III-B now gives this formula and Section~IV-B draws the
consequence: what distinguishes the blue channel is that it is almost empty, with 54.5\% of its
pixels exactly zero, so the trigger is the only strong response there. We state that this saliency
explanation is consistent with the measurements but not established, and name the contrast-matched
trigger as the experiment that would settle it.

\point{The ``serious threat'' framing should be moderated.} The abstract and conclusion now say ``a
serious threat to federated malware classifiers in this controlled IID setting''.

\section{Changes no reviewer asked for}

Two problems surfaced while we prepared the revision, and we report them rather than leave them in.

\point{The seed-42 clean baseline was trained on a different budget.} The clean FL baseline for seed
42, and the trigger controls evaluated against its checkpoint, come from a run of 30 rounds with one
local epoch, while every other run uses 50 rounds with two local epochs. Section~III-H now states
this. We re-trained the baseline under the common budget: it gives 0.7867 accuracy and 0.7869
macro-F1 with a control rate of 6/15 instead of 4/15, so it is worse than the run we report. We state
both directions of the resulting bias: the reported run sets a higher bar for the claim that the
backdoors preserve clean performance, and a lower one for the trigger-control argument. No ASR result
is affected, because every backdoor and defense run uses the common budget.

\point{Table III was not a paired comparison.} The two representations were generated with different
split files, so the RGB-stack and opacity-blend rows were measured on different 75-image test sets
and cannot support a performance claim across representations. Section~III-C now says so, and the
selection of RGB-stack rests on the design argument the paper already made, that channel separation
is required for channel-aware analysis, with macro-F1 selecting only the backbone within that
representation. The abstract, Section~IV-A and the conclusion were adjusted to match.

\point{Figure 7 was regenerated.} The version in the submitted paper was exported at 8.5 inches wide
and then placed in a 3.5 inch column, so its tick labels printed at about 2.3\,pt against 10\,pt body
text. It is now generated at the size it is placed at and its labels print at 6 to 7\,pt. The data is
unchanged: the final-round ASR of all four curves reproduces Table~VI exactly.

\section{What we did not change}

Three of Reviewer C's points need experiments we did not run for this revision: more seeds, a second
source-target pair, and a non-IID partition. All three are now stated as limitations in the
conclusion, together with the note that the red and green results rest on seed 42 alone, and the
three matching experiments are named as future work. We would rather mark the boundary of the
evidence than argue past it.

\end{document}
